% GENERATED FILE. Edit canonical lessons, then run npm run generate:books.
% canonical-source-hash: fnv1a64:a84da3463086b96e
% canonical-lessons: GE-C498-kennenlernen, GE-C498-befreundet, GE-C498-bekanntschaft, GE-C498-nachbarschaft, GE-C498-gastfreundschaft

\chapter{Freunde und Nachbarn}
\label{ch:ge-freunde-und-nachbarn}
\hlchaptermodality{498}

\begin{chapteropening}
\textbf{By the end of this chapter:} \emph{I can talk about friends, neighbours and hospitality.}

Complete the last lesson of chapter 498: I can talk about friends, neighbours and hospitality.
\end{chapteropening}

\section[kennenlernen]{kennenlernen — to get to know, to meet}

\label{lesson:GE-C498-kennenlernen}



\begin{quote}
Before the new one: say the German for to tease, then the German for pocket money.
\end{quote}

\subsection*{You'll want to know: kennenlernen}
\textbf{kennenlernen} — \textquotedblleft{}to get to know, to meet\textquotedblright{}.

\begin{quote}
Wir haben uns im Sprachkurs kennengelernt.
\end{quote}

\begin{quote}
\emph{We met on a language course.}
\end{quote}

Separable: \emph{ich lerne ihn kennen}, \emph{hat kennengelernt}. Also written as two words: \emph{kennen lernen}.

\begin{cousinweb}
\emph{kennenlernen} is built on \emph{kennen} \textquotedblleft{}to know\textquotedblright{}, inherited from Old High German \emph{kennan}, from Proto-Germanic \emph{kannijaną}, and \emph{lernen} \textquotedblleft{}to learn\textquotedblright{} (EtymDB).
\end{cousinweb}

\begin{grammarlens}[title={What you've built}]
One more word for telling, explaining and describing at B1.
\end{grammarlens}

\subsection*{Your turn}
\begin{itemize}
\raggedright
  \item \emph{Say it:} \emph{kennenlernen}
  \item \emph{Say it:} \emph{kennenlernen}, once more
  \item \emph{Say it:} a sentence of your own with \emph{kennenlernen}
  \item \emph{Recall:} say the German for to tease, then the German for pocket money, then say \emph{kennenlernen} again
\end{itemize}

\begin{tcolorbox}[breakable,colback=teal!4,colframe=teal!35!black,title={Before you move on}]
What does kennenlernen mean? (\textquotedblleft{}To get to know, to meet\textquotedblright{}.) Say it once more.
\end{tcolorbox}
\section[befreundet]{befreundet — on friendly terms, friends (with)}

\label{lesson:GE-C498-befreundet}



\begin{quote}
Before the new one: say the German for pocket money, then the German for to get to know.
\end{quote}

\subsection*{You'll want to know: befreundet}
\textbf{befreundet} — \textquotedblleft{}on friendly terms, friends (with)\textquotedblright{}.

\begin{quote}
Ich bin mit meiner Nachbarin seit Jahren befreundet.
\end{quote}

\begin{quote}
\emph{I've been friends with my neighbour for years.}
\end{quote}

\emph{Mit jemandem befreundet sein} — to be friends with someone.

\begin{cousinweb}
\emph{befreundet} is built on \emph{be-} and \emph{Freund} \textquotedblleft{}friend\textquotedblright{}, inherited from Middle High German \emph{vriunt} and Old High German \emph{friunt} (EtymDB).
\end{cousinweb}

\begin{grammarlens}[title={What you've built}]
One more word for telling, explaining and describing at B1.
\end{grammarlens}

\subsection*{Your turn}
\begin{itemize}
\raggedright
  \item \emph{Say it:} \emph{befreundet}
  \item \emph{Say it:} \emph{befreundet}, once more
  \item \emph{Say it:} a sentence of your own with \emph{befreundet}
  \item \emph{Recall:} say the German for pocket money, then the German for to get to know, then say \emph{befreundet} again
\end{itemize}

\begin{tcolorbox}[breakable,colback=teal!4,colframe=teal!35!black,title={Before you move on}]
What does befreundet mean? (\textquotedblleft{}On friendly terms, friends (with)\textquotedblright{}.) Say it once more.
\end{tcolorbox}
\section[die Bekanntschaft]{die Bekanntschaft — acquaintance}

\label{lesson:GE-C498-bekanntschaft}



\begin{quote}
Before the new one: say the German for to get to know, then the German for on friendly terms.
\end{quote}

\subsection*{You'll want to know: die Bekanntschaft}
\textbf{die Bekanntschaft} — \textquotedblleft{}acquaintance\textquotedblright{}.

\begin{quote}
Im Urlaub habe ich eine nette Bekanntschaft gemacht.
\end{quote}

\begin{quote}
\emph{I made a nice acquaintance on holiday.}
\end{quote}

Plural: \emph{die Bekanntschaften}. \emph{Jemandes Bekanntschaft machen} — to get to know someone.

\begin{cousinweb}
\emph{Bekanntschaft} is built on \emph{bekannt} \textquotedblleft{}known\textquotedblright{} with \emph{-schaft}; \emph{bekannt} belongs with \emph{kennen} \textquotedblleft{}to know\textquotedblright{}, from Proto-Germanic \emph{kannijaną} (EtymDB).
\end{cousinweb}

\begin{grammarlens}[title={What you've built}]
One more word for telling, explaining and describing at B1.
\end{grammarlens}

\subsection*{Your turn}
\begin{itemize}
\raggedright
  \item \emph{Say it:} \emph{die Bekanntschaft}
  \item \emph{Say it:} \emph{die Bekanntschaft}, once more
  \item \emph{Say it:} a sentence of your own with \emph{die Bekanntschaft}
  \item \emph{Recall:} say the German for to get to know, then the German for on friendly terms, then say \emph{die Bekanntschaft} again
\end{itemize}

\begin{tcolorbox}[breakable,colback=teal!4,colframe=teal!35!black,title={Before you move on}]
What does die Bekanntschaft mean? (\textquotedblleft{}Acquaintance\textquotedblright{}.) Say it once more.
\end{tcolorbox}
\section[die Nachbarschaft]{die Nachbarschaft — neighbourhood}

\label{lesson:GE-C498-nachbarschaft}



\begin{quote}
Before the new one: say the German for on friendly terms, then the German for acquaintance.
\end{quote}

\subsection*{You'll want to know: die Nachbarschaft}
\textbf{die Nachbarschaft} — \textquotedblleft{}neighbourhood\textquotedblright{}.

\begin{quote}
In unserer Nachbarschaft gibt es einen schönen Park.
\end{quote}

\begin{quote}
\emph{There's a nice park in our neighbourhood.}
\end{quote}

Also the neighbours themselves: \emph{Die ganze Nachbarschaft hat geholfen.}

\begin{cousinweb}
\emph{Nachbarschaft} is built on \emph{Nachbar} \textquotedblleft{}neighbour\textquotedblright{}, inherited from Middle High German \emph{nāchbūre} \textquotedblleft{}one who lives near\textquotedblright{} (EtymDB), with \emph{-schaft}.
\end{cousinweb}

\begin{grammarlens}[title={What you've built}]
One more word for telling, explaining and describing at B1.
\end{grammarlens}

\subsection*{Your turn}
\begin{itemize}
\raggedright
  \item \emph{Say it:} \emph{die Nachbarschaft}
  \item \emph{Say it:} \emph{die Nachbarschaft}, once more
  \item \emph{Say it:} a sentence of your own with \emph{die Nachbarschaft}
  \item \emph{Recall:} say the German for on friendly terms, then the German for acquaintance, then say \emph{die Nachbarschaft} again
\end{itemize}

\begin{tcolorbox}[breakable,colback=teal!4,colframe=teal!35!black,title={Before you move on}]
What does die Nachbarschaft mean? (\textquotedblleft{}Neighbourhood\textquotedblright{}.) Say it once more.
\end{tcolorbox}
\section[die Gastfreundschaft]{die Gastfreundschaft — hospitality}

\label{lesson:GE-C498-gastfreundschaft}



\begin{quote}
Before the new one: say the German for acquaintance, then the German for neighbourhood.
\end{quote}

\subsection*{You'll want to know: die Gastfreundschaft}
\textbf{die Gastfreundschaft} — \textquotedblleft{}hospitality\textquotedblright{}.

\begin{quote}
Vielen Dank für eure Gastfreundschaft!
\end{quote}

\begin{quote}
\emph{Thank you so much for your hospitality!}
\end{quote}

No plural.

\begin{cousinweb}
\emph{Gastfreundschaft} is built on \emph{Gast} \textquotedblleft{}guest\textquotedblright{}, inherited from Old High German \emph{gast}, and \emph{Freundschaft} \textquotedblleft{}friendship\textquotedblright{} (EtymDB).
\end{cousinweb}

\begin{grammarlens}[title={What you've built}]
One more word for telling, explaining and describing at B1.
\end{grammarlens}

\subsection*{Your turn}
\begin{itemize}
\raggedright
  \item \emph{Say it:} \emph{die Gastfreundschaft}
  \item \emph{Say it:} \emph{die Gastfreundschaft}, once more
  \item \emph{Say it:} a sentence of your own with \emph{die Gastfreundschaft}
  \item \emph{Recall:} say the German for acquaintance, then the German for neighbourhood, then say \emph{die Gastfreundschaft} again
\end{itemize}

\begin{tcolorbox}[breakable,colback=teal!4,colframe=teal!35!black,title={Before you move on}]
What does die Gastfreundschaft mean? (\textquotedblleft{}Hospitality\textquotedblright{}.) Say it once more.
\end{tcolorbox}
